% ============================================================================
%  Chapter 7 -- Verm1, Part 1: A Head That Crosses the World     (Lecture 07)
% ============================================================================
\chapter{Verm1, Part 1: A Head That Crosses the World}
\label{ch:verm1p1}

\epigraphline{The journey of a thousand miles begins with a single step.}
  {Laozi, \emph{Tao Te Ching}}

\begin{chapgoals}
By the end of this chapter you will have
\begin{itemize}
\item a cartridge of your own, \fname{verm1.prg32}, built from one assembly
  file that you typed;
\item a title, a grid of 40 by 24 cells, and the head of the worm drawn in one of them;
\item the head moving by itself, one cell at a time, at a speed that you control in fixed point;
\item a world with no edges: past the last column the head comes back from the first;
\item the habit of the working loop: edit, build, run, look, and stop the
  program in the debugger.
\end{itemize}
\end{chapgoals}
% Generated by tools/textbook/export_listings.js from tools/decks/mastery.js and tools/concepts/registry.json: do not edit.
\begin{chapstatus}{Learning status of this chapter}
\masterychip{now}{UNDERSTAND NOW}\enspace the Verm1 project; the init, update, draw contract; a colour is a number; arrays in memory; .text, .rodata, .data; address, value, load, store; why ra is saved; shifts as multiplications; fixed point; pseudo-instructions; little-endian byte order.\par\smallskip
\masterychip{use}{USE NOW}\enspace the PRG32 tools, the PRG32-QT debugger (used, not studied, in this course).\par\smallskip
\masterychip{later}{UNDERSTAND LATER}\enspace instruction formats (Chapter~\ref{ch:instructions}).
\end{chapstatus}


This is the first tutorial of the project. It uses only what Chapters
\ref{ch:firstcontact} to~\ref{ch:reals} have taught, and it is meant to be
followed at the keyboard: the reader types the program. Everything printed in
the listings is read from the reference file \fname{verm1\_part1.S}, which is
built and run automatically; the whole file is in Appendix~\ref{app:verm1}.

% ----------------------------------------------------------------------------
\section{The project}

\subsection{Verm1: a tribute to Snake}

\vermone, read \emph{vermone}, ``the big worm'', is a worm that moves on a
grid, eats, grows, and must not bite itself. The game has a long history: it
appeared as \emph{Blockade}, an arcade machine by Gremlin, in 1976, and as
\emph{Snake} on the Nokia~6110 in 1997, written by Taneli Armanto, it became
one of the most played games ever. \vermone is a tribute, not a copy: the
name, the rules of detail and the look are its own. It is small enough to be
written from scratch in assembly and rich enough to need every topic of the
course, and the reader writes every line of it: there is no engine, and no
library beyond the calls of \prg. The reference implementation that the
tutorial chapters follow, and against which a reader can always compare a
program of their own, was contributed by Giuseppe Coviello
(Appendix~\ref{app:verm1}). Figure~\ref{fig:v1:goal} shows the destination.

\begin{figure}[htb]
\centering
\includegraphics[width=0.55\textwidth]{verm1_p4_play.png}
\caption{\vermone at the end of the book (Part 4).}\label{fig:v1:goal}
\end{figure}

\subsection{The plan}

\begin{table}[htb]
\centering\small
\caption{Four parts and a publication.}\label{tab:v1:plan}
\begin{tabularx}{\textwidth}{@{}lLl@{}}
\toprule
\textbf{Chapter} & \textbf{You add} & \textbf{It uses Chapters}\\
\midrule
\ref{ch:verm1p1}\quad Part 1 & a grid, a head, automatic movement, wrap-around & \ref{ch:history} to \ref{ch:reals}: information and numbers\\
\ref{ch:verm1p2}\quad Part 2 & steering with the joystick, food, a score & \ref{ch:boole1} to \ref{ch:combinational}: logic and circuits\\
\ref{ch:verm1p3}\quad Part 3 & a body that grows, collisions, title and game-over screens & \ref{ch:sequential} to \ref{ch:storage}: state and memory\\
\ref{ch:verm1p4}\quad Part 4 & fast drawing, a score in digits, levels, sound & \ref{ch:processor} to \ref{ch:io}: the machine\\
\ref{ch:store}\quad The Store & a bundle with manifest, icon and credits & all of them\\
\bottomrule
\end{tabularx}
\end{table}

Each part ends with a program that runs (Table~\ref{tab:v1:plan}).

\begin{prgwarn}
Never start a part with a program that does not run. If yours is broken,
start from the reference file of the previous part.
\end{prgwarn}

Part~1 takes the following from the chapters read so far.

\begin{center}\small
\begin{tabularx}{\textwidth}{@{}lLL@{}}
\toprule
\textbf{From} & \textbf{Idea} & \textbf{Where it appears}\\
\midrule
Chapter \ref{ch:firstcontact} & registers, \ins{la}/\ins{lw}/\ins{sw}, \ins{call} and the stack frame & everywhere\\
Chapter \ref{ch:history} & retro games as specimens & the choice of the game\\
Chapter \ref{ch:information} & colours as RGB565 numbers; strings ending with zero & \code{GREEN}, the title\\
Chapter \ref{ch:integers} & shifts as multiplications; comparisons; wrap-around & cell to pixel; the edge of the world\\
Chapter \ref{ch:reals} & fixed point: a fraction as a scaled integer & the speed of the worm\\
\bottomrule
\end{tabularx}
\end{center}

% ----------------------------------------------------------------------------
\section{Step 0: the working loop}

Make a folder \fname{\textasciitilde/verm1} and keep in it one file,
\fname{verm1.S}, which will grow from part to part. The two commands of
Listing~\ref{lst:v1:build} build the cartridge and send it to \prgqt; they
are the same in every tutorial. \prgqt listens on port 8080 and shows its
address on the Setup screen: use that address if \code{localhost} does not
answer. The same cartridge runs on QEMU (\code{qemu upload}, then
\code{qemu run}) and on the board (\code{-{}-url http://192.168.4.1}).

\shfile[caption={Build and run: edit, build, run, look.},label=lst:v1:build]{code/l07/c03.sh}

% ----------------------------------------------------------------------------
\section{Step 1: three functions that return}

\asmfile[caption={Step 1 \textbf{(new)}: the skeleton.},label=lst:v1:skel]{code/l07/c04.S}

This is the skeleton of Chapter~\ref{ch:firstcontact} with another prefix:
\code{verm1} is the name passed to \code{-{}-entry-prefix}. \code{init} runs
once; \code{update} and \code{draw} run once per frame.

\begin{prgcheck}[Run it]
Type it, build it, run it: a black screen is the right answer. If the build
fails here the problem is the installation, not the program: solve it before
going on. From now on, add one thing, run, look. Never two things at once.
\end{prgcheck}

% ----------------------------------------------------------------------------
\section{Step 2: a grid and a head}

\subsection{Constants: colours and the grid}

Replace the first lines of the skeleton with those of
Listing~\ref{lst:v1:const}.

\asmfile[caption={Step 2 \textbf{(changed)}: the head of the file.},label=lst:v1:const]{code/l07/c05.S}

Colours are RGB565 numbers: \hex{07e0} is pure green. The screen is 320
pixels across and a cell is 8 pixels wide, which makes 40 columns. The top 8
pixels hold the title, and 192 are left: 24 rows. The directive \code{.equ}
gives a name to a number and uses no memory. The constant \code{SPEED\_Q8} is
used in step~3.

\subsection{The world of Verm1}

\begin{prgconcept}
The worm does not live in pixels: it lives in cells. A position is a column,
0 to 39, and a row, 0 to 23.
\end{prgconcept}

To draw a cell its coordinates are turned into pixels
(Figure~\ref{fig:v1:grid}): $x = \mathit{column}\times 8$, a shift to the left
by 3, and $y = \mathit{row}\times 8+8$, because the title line comes first.
The cell $(0,0)$ is at pixel $(0,8)$; the cell $(39,23)$ is at $(312,192)$.
Thinking in cells makes every later rule simple: moving is adding~1.

\begin{figure}[htb]
\centering
\begin{tikzpicture}[x=2.6mm,y=-2.6mm]
  \fill[prgsteel] (0,0) rectangle (40,1);
  \fill[prgink] (0,1) rectangle (40,25);
  \draw[prgsteel,very thin,step=1] (0,1) grid (40,25);
  \fill[green!80!black] (5,13) rectangle (6,14);
  \node[text=white,font=\sffamily\tiny,anchor=west] at (0.3,0.5) {title line: 8 pixels};
  \node[lbl] at (20,-1) {columns $0\ldots 39$\quad ($x=\mathit{column}\times 8$)};
  \node[lbl,rotate=90] at (-3.2,13) {rows $0\ldots 23$\quad ($y=\mathit{row}\times 8+8$)};
  \node[lbl] at (5.5,26) {column 5};
  \node[lbl,anchor=east] at (-0.3,13.5) {row 12};
\end{tikzpicture}
\caption{The world of \vermone: 40 by 24 cells of 8 by 8 pixels. The head
starts in column 5 of row 12.}
\label{fig:v1:grid}
\end{figure}

\subsection{Variables and a string}

Put Listing~\ref{lst:v1:data} at the end of the file. The section
\code{.rodata} holds constants: the title ends with a zero byte. The section
\code{.data} holds variables, one 32-bit word each. Their initial values are
only a start: \code{init} will set them again.

\asmfile[caption={Step 2 \textbf{(new)}: data, at the end of the file.},label=lst:v1:data]{code/l07/c06.S}

\subsection{\texttt{init}: put the head in its place}

\asmfile[caption={Step 2 \textbf{(changed)}: \code{verm1\_init}.},label=lst:v1:init]{code/l07/c07.S}

\ins{la} takes an \emph{address}, \ins{sw} stores a \emph{value} there. The
function sets the start explicitly: a program should never trust what memory
contained. \reg{zero} is the register that always reads 0. The function is a
leaf: it makes no calls, so it needs no stack frame.

\subsection{\texttt{draw}: title and head}

\asmfile[caption={Step 2 \textbf{(changed)}: \code{verm1\_draw}.},label=lst:v1:draw]{code/l07/c08.S}

Clear the screen, write the title, draw one rectangle: three calls, with the
arguments travelling in \reg{a0} to \reg{a4}, in order. The signatures are
\code{prg32\_gfx\_text8(x, y, string, foreground, background)} and
\code{prg32\_gfx\_rect(x, y, width, height, colour)}. \ins{slli} by 3
multiplies by 8: from cell to pixel. Since \code{draw} makes calls, it saves
\reg{ra} in a stack frame.

\begin{prgcheck}[Run it: a head on the grid]
Expected now: the word \texttt{VERM1} at the top left and a green square
standing still in column 5, row 12 (Figure~\ref{fig:v1:shot} shows it later,
when it has begun to move).
\end{prgcheck}

\begin{figure}[htb]
\centering
\includegraphics[width=0.55\textwidth]{verm1_p1.png}
\caption{The title and the head. This picture was taken by the automatic
tests 64 frames into the finished Part~1: the head is in column 15.}
\label{fig:v1:shot}
\end{figure}

% ----------------------------------------------------------------------------
\section{Step 3: movement}

\subsection{How fast should a worm go?}

\code{update} runs about thirty times per second. One cell per frame would
cross the screen in little more than a second: unplayable. What is wanted is
a fraction of a cell per frame, and the processor of the kit has only
integers. The answer is the fixed point of Chapter~\ref{ch:reals}, in the
format Q24.8: the program keeps a fraction in units of $1/256$ of a cell, and
takes a step each time a whole cell has accumulated. With
$\code{SPEED\_Q8}=40$ the worm advances $40/256\approx 0.156$ cells per frame:
about one step every 6 frames.

This is the ball of Chapter~\ref{ch:reals} turned into a timer. The remainder
is not thrown away, so the average speed is exact: 256 frames give exactly 40
steps.

\subsection{\texttt{update}: accumulate, step, wrap}

\asmfile[caption={Step 3 \textbf{(changed)}: \code{verm1\_update}.},label=lst:v1:update]{code/l07/c10.S}

The speed is added to the fraction; below 256, the function waits. At 256 it
subtracts 256, keeps the rest, and moves the head one cell to the right.
Column 40 does not exist: past the right edge the head goes back to column~0.
The function has two exits, and each stores what it has changed. \ins{blt}
compares signed numbers, which is right here because every value is small and
positive. The local labels \code{1:} and \code{2:} are referred to as
\code{1f} and \code{2f}: ``the next 1, or 2, \emph{forward}''.

\subsection{Trace it by hand}

\begin{table}[htb]
\centering\small
\caption{A trace of \code{verm1\_update}.}\label{tab:v1:trace}
\begin{tabular}{@{}cccccc@{}}
\toprule
\textbf{Frame} & \textbf{\code{step\_q8} before} & \textbf{+ 40} & \textbf{Whole cell?} & \textbf{\code{step\_q8} after} & \textbf{\code{head\_x}}\\
\midrule
1 & 0 & 40 & no & 40 & 5\\
2 & 40 & 80 & no & 80 & 5\\
6 & 200 & 240 & no & 240 & 5\\
7 & 240 & 280 & yes: $280-256$ & 24 & 6\\
13 & 224 & 264 & yes: $264-256$ & 8 & 7\\
64 & -- & -- & -- & 0 & 15\\
\bottomrule
\end{tabular}
\end{table}

The steps come every 6 or 7 frames (Table~\ref{tab:v1:trace}): the remainder
carries the difference, as a leap year does. After 64 frames the sum is
$64\times 40 = 2560 = 10\times 256$: ten steps, and the head is in column 15,
exactly as in Figure~\ref{fig:v1:shot}. Fill in frames 8 to 12 as an
exercise.

\begin{prgcheck}[Run it: the head crosses the world]
\begin{enumerate}[nosep]
\item Build and run: the head moves to the right, about five cells per second.
\item Wait for the right edge: the head leaves at column 39 and returns at
  column 0.
\item Change \code{SPEED\_Q8} to 8, then to 128, then to 256: predict the
  speed before you run.
\item Change the starting row in \code{init}, and the colour of the head.
\end{enumerate}
\textbf{Expected result.} The head crosses the forty columns in 256 frames,
about eight and a half seconds at thirty frames per second, for ever. With
$\code{SPEED\_Q8}=256$ it moves one cell per frame; larger values still give
one step per frame, because \code{update} takes at most one step: a limit
worth discovering.
\end{prgcheck}

% ----------------------------------------------------------------------------
\section{Step 4: look inside with PRG32-QT}

The debugger is part of the method, not an emergency tool.

\begin{prgtry}[Stop the program and read the machine]
\begin{enumerate}[nosep]
\item In \prgqt open \emph{Settings} and enable \emph{Run cartridges in RISC-V
  debug mode}.
\item Run your cartridge: the game is on the left, the debugger on the right.
\item Press \key{Update}: the listing jumps to \code{verm1\_update}. Click in
  the gutter to put a breakpoint on its first instruction.
\item The program stops there at the next frame. Press \key{Step} and watch
  the registers change, one instruction at a time.
\item After the two instructions of \ins{la}, \reg{t0} holds an address: type
  it in the \emph{Memory} field to see \code{step\_q8} in memory.
\item Press \key{Resume}, and watch the bytes of \code{step\_q8} grow by 40
  (\hex{28}) at every stop.
\end{enumerate}
\end{prgtry}

\key{Resume} ignores the current breakpoint once, so that each press runs one
whole frame.

\subsection{What the listing shows}

The debugger does not show the source file: it shows the code as the
assembler left it (Listing~\ref{lst:v1:dis}).

\txtfile[caption={\code{verm1\_update}, disassembled (abridged; the exact offsets depend on the build).},label=lst:v1:dis]{code/l07/c12.txt}

\begin{itemize}
\item Names are gone: labels have become addresses, and constants have become
  numbers.
\item Pseudo-instructions are expanded: \ins{la} is \ins{auipc} plus
  \ins{addi}, \ins{li} is \ins{addi} from \reg{zero}, \ins{ret} is a jump to
  \reg{ra}.
\item Some instructions start with \code{c.}: they are the \emph{compressed},
  16-bit forms of the same instructions (the C of RV32IMAC).
\item The cartridge is loaded at address \hex{40800000}: every address shown
  starts there.
\end{itemize}

Instruction formats are the subject of Chapter~\ref{ch:instructions}; the
compressed forms are only named there, and are listed in Appendix~\ref{app:riscv}.

\subsection{Reading memory: little-endian, live}

The \emph{Memory} view shows 128 bytes from the address typed, in hexadecimal
(Figure~\ref{fig:v1:debug}). The variable \code{step\_q8} is a word: its four
bytes come lowest byte first, so that with $\code{step\_q8}=240$ one reads
\texttt{F0 00 00 00}. In the file the order of the variables is
\code{head\_x}, \code{head\_y}, \code{step\_q8}, four bytes each: if
\code{step\_q8} is at address $A$, \code{head\_x} is at $A-8$.

The debugger only observes. To change the program: edit, build, run.

\begin{figure}[htb]
\centering
\includegraphics[width=0.9\textwidth]{prg32qt_debugger_details.png}
\caption{\prgqt: listing, registers, memory.}\label{fig:v1:debug}
\end{figure}

% ----------------------------------------------------------------------------
% Generated by tools/textbook/export_listings.js from tools/concepts/assessments.json: do not edit.
\begin{chapstatus}{The thread of this chapter}
Numbers became positions, colours and a speed: the whole information block, in one running game. Next: the logic that will let the player steer it.
\par\medskip\textbf{Exit ticket} (two to four minutes, alone, on paper):
\begin{enumerate}[nosep,leftmargin=2em]
\item Verm1 declares title: .asciz "VERM1" and head\_x: .word 5. In which section does each belong, and why must head\_x not be placed in .rodata?
\item step\_q8 grows by SPEED\_Q8 at every frame, and the head moves one cell each time it reaches 256. With SPEED\_Q8 = 64, how many frames between two steps? With 128? Predict, then run.
\end{enumerate}
\end{chapstatus}


\section{Check your understanding}

\begin{quiz}{With $\code{SPEED\_Q8}=64$, how many frames pass between two steps of the head?}
\item 2
\item 4
\item 64
\item 256
\end{quiz}

\begin{quiz}{Which pixel is the top-left corner of the cell in column 10, row 3?}
\item $(10, 3)$
\item $(80, 24)$
\item $(80, 32)$
\item $(88, 32)$
\end{quiz}

\begin{quiz}{Why does \code{verm1\_draw} save \reg{ra} while \code{verm1\_update} does not?}
\item \code{draw} is longer
\item \code{draw} calls other functions, which overwrite \reg{ra}
\item \code{update} uses no registers
\item \reg{ra} is saved automatically in \code{update}
\end{quiz}

\section{Make it yours}

\begin{exercise}[on \prg]
\begin{enumerate}[nosep]
\item Make the head move to the left: which instruction changes, and which
  comparison detects the edge?
\item Make it move downwards instead, wrapping from row 23 to row 0.
\item Draw a second cell, in a darker green, one column behind the head: a
  first hint of a body.
\item Write your own name next to the title, in a colour that you computed by
  hand in RGB565.
\end{enumerate}
\expected{Each change is a few lines. Keep a copy of the working Part~1
before experimenting: Part~2 starts from it. (Hints: to move left,
\code{addi t3, t3, -1}, then \code{bgez t3} to skip the wrap, else
\code{li t3, COLS-1}: the comparison is on the sign. A darker green is
\hex{03e0}.)}
\end{exercise}

\section{Where we are, and what comes next}

The reader has a cartridge that runs, built from a file typed and understood
line by line. It has state in memory, a rule that changes the state at every
frame, and a function that draws it: the shape of every game. What is missing
is the player: the head goes where it wants. Chapters \ref{ch:boole1}
to~\ref{ch:combinational} give the tools: bits and masks to read the joystick,
Boolean rules, tables that select. Part~2 adds steering, food and a score.
